\documentclass[11pt]{article}
\usepackage{fontspec,amsmath,amssymb,amsthm,hyperref}
\usepackage[margin=1in]{geometry}
\setlength{\emergencystretch}{3em}
\title{The actual mean-field force is Lipschitz in W\ensuremath{{}_{2}}}
\author{LangevinLean mathematical exposition}
\date{}
\begin{document}\maketitle
\noindent Source module: \texttt{ProximalLangevin/MeanFieldDrift.lean}. Compile this proof document with XeLaTeX or LuaLaTeX. The accompanying HTML places the complete Lean declarations beside every proof.


For a pair potential \(V\) with \(L\)-Lipschitz gradient, define \(b_V(\theta,\mu)=\int\nabla_\theta V(\theta,x)\mu(dx)\). The measures are probability laws with finite second moments. Every Bochner expectation and every finite real transport distance is justified.


\section{A Lipschitz observable preserves L\ensuremath{{}^{2}} and controls expectation differences}


Center an observable \(\phi\) at \(\phi(0)\); its norm growth is bounded by \(K\|x\|\), so \(\phi(X)\in L^2\subset L^1\). Given any finite-cost coupling \(c\) of \(\mu,\nu\), pushforward integration and Cauchy–Schwarz give

\[\left\|\int\phi\,d\mu-\int\phi\,d\nu\right\|\le K\left(\int\|x-y\|^2dc\right)^{1/2}.\]

Taking the actual coupling infimum yields the bound \(K W_2(\mu,\nu)\), with finiteness supplied by the second moments.



\noindent Lean declarations: \texttt{memLp\_lipschitz\_observable}, \texttt{integral\_lipschitz\_wasserstein2}.

\section{Project the true joint gradient to its parameter block}


The bounded linear map \(P_\theta\) keeps the parameter coordinates. Its squared norm is a sub-sum of the full Euclidean square, hence \(\|P_\theta z\|\le\|z\|\). Define \(\phi_\theta(x)=P_\theta\nabla V(\theta,x)\), then define the mean-field force as its actual Bochner integral. These definitions use the existing joint potential, with no arbitrary force field substituted for its gradient.



\noindent Lean declarations: \texttt{jointParameterCLM}, \texttt{jointParameterCLM\_apply}, \texttt{norm\_jointParameterCLM\_apply\_le}, \texttt{parameterGradient}, \texttt{meanFieldDrift}.

\section{Control joint parameter and latent perturbations}


Decompose a joint pair difference into its two coordinate inclusions, giving \(\|(\theta,x)-(\theta',x')\|\le\|\theta-\theta'\|+\|x-x'\|\). Gradient Lipschitzness and the nonexpansive parameter projection imply

\[\|\phi_\theta(x)-\phi_{\theta'}(x')\|\le L(\|\theta-\theta'\|+\|x-x'\|).\]

Fixing \(\theta=\theta'\) proves that the latent observable itself is \(L\)-Lipschitz.



\noindent Lean declarations: \texttt{pair\_norm\_sub\_le}, \texttt{parameterGradient\_difference\_le}, \texttt{parameterGradient\_latent\_lipschitz}.

\section{Separate the parameter change from the law change}


For a differentiable potential \(V\), define


\[
b_V(\theta,\mu)=\int \nabla_\theta V(\theta,x)\,\mu(dx).
\]


Assume \(\nabla V\) is \(L\)-Lipschitz and \(\mu,\nu\) are probability laws with finite second moment. Coordinate projection has norm at most one, so the integrand is Lipschitz and grows at most linearly. Its Bochner expectation is therefore integrable. For any finite-cost coupling, compare its two integrands, apply the norm-of-integral inequality, and use Cauchy–Schwarz. Taking the infimum of actual coupling costs gives


\[
\|b_V(\theta,\mu)-b_V(\theta',\nu)\|
 \le L\bigl(\|\theta-\theta'\|+W_2(\mu,\nu)\bigr).
\]


With the genuine split envelope \(V=g_1+g_2^\eta\), the proved constant is \(L_1+1/\eta\). The paper’s \(1/\eta\) constant applies to a full envelope, not to this split expression in general. This is a corrected \textbf{W\ensuremath{{}_{2}} version} of Lemma C.1: its right-hand side is weaker than the paper’s W\ensuremath{{}_{1}} statement, and the distinction is explicit. Propositions C.2–C.3 are not consequences formalized by this lemma alone.



\noindent Lean declarations: \texttt{meanFieldDrift\_lipschitz\_wasserstein2}, \texttt{moreau\_meanFieldDrift\_lipschitz\_wasserstein2}.

\end{document}
